\section{Comparison with other codes and studies}
\label{sec:code_comparison}

Table~\ref{tab:lit_comparison} summarizes literature-reported deposition strategies alongside metrics from this program.
External entries are \emph{not} direct binary benchmarks against our executable: dimensionality, physics module (electrostatic charge versus electromagnetic current), hardware, and reported metrics differ by design.
The table positions our results relative to published engineering conclusions.

\begin{table}[t]
\centering
\small
\caption{Literature comparison of deposition engineering (metrics are not directly commensurate).}
\label{tab:lit_comparison}
\resizebox{\linewidth}{!}{%
\begin{tabular}{@{}llll@{}}
\toprule
Code / study & Focus & Literature metric & This program \\
\midrule
Harvey et al.\ \cite{harvey2015} & Hybrid charge-conserving deposition & $\sim$4$\times$ GPU vs CPU baseline & CIC GPU atomics 5.3$\times$ vs T4-host CPU \\
PIConGPU \cite{burau2010,steiniger2023} & 3D EM, Esirkepov/EZ, supercell cache & Qualitative GPU cache strategy & GPU atomics 0.81\,ms deposit \\
Smilei \cite{beck2019} & Cycle sort + vectorized projection & $\sim$2--3$\times$ operator speedup & Sort 2.3\,ms vs deposit 0.59\,ms \\
Viereck CPC \cite{viereck2016} & Architecture survey & Memory-bound scatter; tiled privatization & Figs.~\ref{fig:roofline},~\ref{fig:bandwidth} \\
\textbf{This program} & 2D electrostatic reference & Deposit ms, $\epsilon_Q$, noise, $\omega$, $\gamma$ & Archived CSV suite \\
\bottomrule
\end{tabular}%
}
\end{table}

Harvey et al.\ demonstrate that tiled privatization can outperform naive GPU scatter for charge-conserving deposition on hybrid nodes; our GPU privatized tiles preserve conservation but remain slower than atomics because each spatial tile scans all particles (Appendix~\ref{app:gpu}).
PIConGPU and Smilei integrate deposition into full 3D electromagnetic timestep loops with production I/O and MPI domain decomposition; our program isolates the deposition kernel and couples it to a minimal FFT Poisson solver for controlled experiments.
Smilei's cycle-sort plus vectorization strategy aligns with our finding that sorting cost must be reported separately from deposit-kernel time when evaluating locality optimizations.
